% =============================================================================
% APPENDIX AAA: CORE-WHO: The Welfare Hierarchy
% =============================================================================
% Module: BCM2_03_SRL (Social Relation Levels)
% Category: CORE
% Version: 1.0 (January 2026)
% Status: Final
% Language: English
%
% PURPOSE: This appendix answers the fundamental question "Who can have utility?"
% It formalizes the Welfare Hierarchy, establishing that utility can exist at
% multiple levels: Individual (L=0), Dyadic (L=1), Group (L=2), ..., Society (L=n).
% The hierarchy is not merely descriptive but prescriptive for welfare aggregation.
%
% 10C POSITION: This is CORE 1 of 10 - the foundational "WHO" question.
%
% =============================================================================

% -----------------------------------------------------------------------------
% HEADER BLOCK (Required)
% -----------------------------------------------------------------------------

\begin{tcolorbox}[colback=gray!5!white, colframe=gray!75!black,
    title=Appendix Header: AAA CORE-WHO]

\begin{tabular}{@{}ll@{}}
\textbf{Code:} & AAA \\
\textbf{Category:} & CORE (Fundamental Theory) \\
\textbf{Full Name:} & CORE-WHO: The Welfare Hierarchy \\
\textbf{Module:} & BCM2\_03\_SRL (Social Relation Levels) \\
\textbf{Version:} & 1.0 (January 2026) \\
\textbf{Status:} & Final \\
\textbf{Language:} & English \\
\end{tabular}

\vspace{0.5em}
\textbf{Purpose:} Formalizes the \textbf{Welfare Hierarchy}, answering the
first fundamental question of the EBF: \textit{``Who can have utility?''}
Establishes that utility exists at multiple aggregation levels $L \in \{0, 1, 2, \ldots, n\}$,
from individual to society, with formal rules for cross-level aggregation.

\textbf{10C Position:} This is CORE 1 (WHO) --- the foundation upon which
WHAT (C), HOW (B), and WHEN (V) build.

\end{tcolorbox}

% =============================================================================
% CROSS-REFERENCE STRUCTURE
% =============================================================================
%
% CHAPTER LINKAGE (Required):
% - Primary Chapter: Chapter 5: The Welfare Hierarchy
% - Secondary Chapters: Ch. 3 (Utility Foundations), Ch. 6 (Social Preferences)
%
% DEPENDENCIES (what this appendix REQUIRES):
% - None (foundational - other COREs build on this)
%
% DEPENDENTS (what REQUIRES this appendix):
% - Appendix C (CORE-WHAT): Needs WHO to specify whose utility dimensions
% - Appendix B (CORE-HOW): Needs WHO for inter-level complementarity
% - Appendix V (CORE-WHEN): Needs WHO for context-sensitive aggregation
% - Appendix AU (CORE-AWARE): Awareness exists at each level L
% - Appendix AV (CORE-READY): Willingness exists at each level L
% - Appendix AW (CORE-STAGE): Journey stages apply per level
% - Appendix HI (CORE-HIERARCHY): Decision hierarchy presupposes welfare hierarchy
% - Appendix IE (CORE-EIT): Interventions target specific levels
% - Appendix CAL (METHOD-LLMMC): Calibration methods for level-specific parameters
%
% =============================================================================

% -----------------------------------------------------------------------------
% CHAPTER-APPENDIX LINKAGE BOX
% -----------------------------------------------------------------------------

\begin{tcolorbox}[colback=purple!5!white, colframe=purple!75!black,
    title=Chapter Linkage]

\textbf{Primary Chapter:} Chapter 5: The Welfare Hierarchy

\vspace{0.3em}
\textbf{The Fundamental Question:}
\begin{quote}
\textit{``Who can have utility? At what levels does welfare exist, and how do we aggregate across levels?''}
\end{quote}

\textbf{Section-Level Mapping:}
\begin{itemize}[nosep]
\item Section 5.1: Why "Who" Matters $\leftrightarrow$ Section~\ref{sec:who-fundamental}
\item Section 5.2: The Level Hierarchy $\leftrightarrow$ Section~\ref{sec:who-levels}
\item Section 5.3: Aggregation Rules $\leftrightarrow$ Section~\ref{sec:who-aggregation}
\item Section 5.4: Cross-Level Externalities $\leftrightarrow$ Section~\ref{sec:who-externalities}
\end{itemize}

\textbf{Secondary Chapters:}
\begin{itemize}[nosep]
\item Chapter 3: Utility Foundations --- Individual utility as base case
\item Chapter 6: Social Preferences --- How $L > 0$ preferences form
\item Chapter 11: Awareness --- Aggregation axioms AWX-A76 to A80
\end{itemize}

\textbf{Relationship to Other COREs:}
\begin{itemize}[nosep]
\item This appendix (AAA): \textbf{WHO} --- Defines the levels at which utility exists
\item Appendix C (CORE-WHAT): \textbf{WHAT} --- Defines the dimensions of utility at each level
\item Appendix B (CORE-HOW): \textbf{HOW} --- Defines how utilities interact across levels
\item Appendix V (CORE-WHEN): \textbf{WHEN} --- Defines when context modifies level weights
\end{itemize}

\end{tcolorbox}

% -----------------------------------------------------------------------------
% APPENDIX SCOPE BOX
% -----------------------------------------------------------------------------

\begin{tcolorbox}[
    colback=orange!5!white,
    colframe=orange!75!black,
    title={\textbf{Appendix Scope: Ziel / In-Scope / Out-of-Scope / Constraints}},
    fonttitle=\bfseries
]

\textbf{Ziel (Objective):}
\begin{itemize}[nosep]
    \item Define the complete hierarchy of utility bearers from individual to society
\end{itemize}

\textbf{In-Scope:}
\begin{itemize}[nosep]
    \item Level definitions ($L_0$ through $L_n$)
    \item Aggregation axioms (SRL-1 through SRL-8)
    \item Cross-level externality formalization
    \item Welfare weights $\omega_L$ specification
\end{itemize}

\textbf{Out-of-Scope (delegiert an andere Appendices):}
\begin{itemize}[nosep]
    \item Utility dimension content $\rightarrow$ Appendix C (CORE-WHAT)
    \item Inter-level complementarity mechanics $\rightarrow$ Appendix B (CORE-HOW)
    \item Context-dependent weighting $\rightarrow$ Appendix V (CORE-WHEN)
    \item Parameter estimation for $\omega_L$ $\rightarrow$ Appendix BBB (CORE-WHERE)
\end{itemize}

\textbf{Constraints:}
\begin{itemize}[nosep]
    \item Levels must be mutually exclusive and collectively exhaustive
    \item Aggregation must satisfy Pareto consistency (SRL-6)
    \item No utility without a bearer (SRL-1)
\end{itemize}

\textbf{Lieferobjekte (Deliverables):}
\begin{itemize}[nosep]
    \item Complete Level Taxonomy (Table~\ref{tab:who-levels})
    \item Axiom System SRL-1 through SRL-8
    \item Aggregation Formula (Equation~\ref{eq:who-aggregation})
    \item Integration specifications with C, B, V
\end{itemize}

\end{tcolorbox}

% =============================================================================
% CHAPTER TITLE
% =============================================================================

\chapter{CORE-WHO: The Welfare Hierarchy}
\label{app:core-who}

\begin{abstract}
This appendix answers the first fundamental question of the Evidence-Based Framework:
\textit{``Who can have utility?''} We establish that welfare exists at multiple
aggregation levels $L \in \{0, 1, 2, \ldots, n\}$, from individual ($L=0$) through
dyadic ($L=1$), group ($L=2$), organizational ($L=3$), to societal ($L=n$) levels.
The axiom system SRL-1 through SRL-8 formalizes the conditions for well-defined
cross-level aggregation, ensuring Pareto consistency and addressing the problem
of double-counting. This hierarchy is foundational: WHAT (C), HOW (B), and WHEN (V)
all presuppose answers to WHO.
\end{abstract}

% -----------------------------------------------------------------------------
% QUICK REFERENCE BOX
% -----------------------------------------------------------------------------

\begin{tcolorbox}[colback=blue!5!white, colframe=blue!75!black,
    title=Quick Reference: CORE-WHO]

\textbf{Core Question:} Who can have utility? At what levels does welfare exist?

\textbf{Key Formula:}
\begin{equation}
W = \sum_{L=0}^{n} \omega_L \cdot U_L \quad \text{where} \quad \sum_{L=0}^{n} \omega_L = 1
\label{eq:who-aggregation}
\end{equation}

\textbf{Key Concepts:}
\begin{itemize}[nosep]
\item \textbf{Level $L$}: Aggregation tier (0=individual, 1=dyad, 2=group, ...)
\item \textbf{Welfare Weight $\omega_L$}: Normative importance assigned to level $L$
\item \textbf{Cross-Level Externality}: Impact of one level's utility on another
\item \textbf{Aggregation}: Combining utilities across levels into total welfare $W$
\end{itemize}

\textbf{Cross-References:}
C (CORE-WHAT), B (CORE-HOW), V (CORE-WHEN), AU (CORE-AWARE), CAL (METHOD-LLMMC)

\end{tcolorbox}


%%%%%%%%%%%%%%%%%%%%%%%%%%%%%%%%%%%%%%%%%%%%%%%%%%%%%%%%%%%%%%%%%%%%%%%%%%%%%%%
\section{The Fundamental Question}
\label{sec:who-fundamental}
%%%%%%%%%%%%%%%%%%%%%%%%%%%%%%%%%%%%%%%%%%%%%%%%%%%%%%%%%%%%%%%%%%%%%%%%%%%%%%%

\subsection{Why "Who" Matters}

Before we can analyze \textit{what} utility consists of, \textit{how} components interact,
or \textit{when} context matters, we must first establish \textit{who} can have utility.
This is not a trivial question:

\begin{itemize}
\item Can a corporation have utility independent of its employees?
\item Does a family's welfare reduce to the sum of individual utilities?
\item Can future generations be utility bearers in present decisions?
\item Does a community have interests beyond those of its members?
\end{itemize}

Different answers lead to radically different policy implications.

\subsection{The Core Problem}

Standard economics typically assumes utility exists only at the individual level ($L=0$),
with social welfare as a mere aggregation function. This creates two problems:

\begin{enumerate}
\item \textbf{Ontological Reductionism}: It cannot capture genuinely collective goods
      (team spirit, social trust, cultural identity) that are not decomposable.
\item \textbf{Aggregation Paradox}: Simple summation ignores that individuals care
      about \textit{distributions} and \textit{relationships}, not just totals.
\end{enumerate}

\begin{tcolorbox}[colback=red!5!white, colframe=red!75!black,
    title=The Fundamental Question]
\textbf{Who can have utility? At what levels does welfare exist, and how do these levels relate to each other?}
\end{tcolorbox}


%%%%%%%%%%%%%%%%%%%%%%%%%%%%%%%%%%%%%%%%%%%%%%%%%%%%%%%%%%%%%%%%%%%%%%%%%%%%%%%
\section{Core Theory: The Level Hierarchy}
\label{sec:who-levels}
%%%%%%%%%%%%%%%%%%%%%%%%%%%%%%%%%%%%%%%%%%%%%%%%%%%%%%%%%%%%%%%%%%%%%%%%%%%%%%%

\subsection{Definition of Levels}

We define the welfare hierarchy as a discrete set of aggregation levels:

\begin{definition}[Welfare Level]
A \textbf{welfare level} $L \in \{0, 1, 2, \ldots, n\}$ is a tier of social organization
at which utility can be meaningfully defined and measured.
\end{definition}

\begin{table}[htbp]
\centering
\caption{The Welfare Level Hierarchy}
\label{tab:who-levels}
\renewcommand{\arraystretch}{1.3}
\begin{tabular}{@{}clll@{}}
\toprule
\textbf{Level} & \textbf{Name} & \textbf{Unit} & \textbf{Example Utility Bearers} \\
\midrule
$L_0$ & Individual & Person & Consumers, voters, employees \\
$L_1$ & Dyadic & Pair & Couples, business partners, mentor-mentee \\
$L_2$ & Group & Team & Families, work teams, friend circles \\
$L_3$ & Organizational & Firm/Institution & Companies, universities, NGOs \\
$L_4$ & Community & Network & Neighborhoods, professional communities \\
$L_5$ & Regional & Territory & Cities, cantons, states \\
$L_n$ & Societal & Society & Nations, global community \\
\bottomrule
\end{tabular}
\end{table}

\subsection{Properties of the Hierarchy}

The hierarchy satisfies three key properties:

\begin{enumerate}
\item \textbf{Nesting}: Each level $L > 0$ contains members from level $L-1$
\item \textbf{Emergence}: Higher-level utility is not fully reducible to lower levels
\item \textbf{Interaction}: Levels influence each other through externalities
\end{enumerate}

\subsection{Formal Structure}

Let $\mathcal{E}_L$ denote the set of entities at level $L$. For entity $e \in \mathcal{E}_L$:

\begin{equation}
U_e^{(L)} = f_L\left( \{U_m^{(L-1)}\}_{m \in e}, \; \Xi_L \right)
\end{equation}

where:
\begin{itemize}
\item $\{U_m^{(L-1)}\}_{m \in e}$ are utilities of constituent members
\item $\Xi_L$ captures emergent properties not reducible to members
\item $f_L$ is the level-specific aggregation function
\end{itemize}


%%%%%%%%%%%%%%%%%%%%%%%%%%%%%%%%%%%%%%%%%%%%%%%%%%%%%%%%%%%%%%%%%%%%%%%%%%%%%%%
\section{Axioms and Formal Definitions}
\label{sec:who-axioms}
%%%%%%%%%%%%%%%%%%%%%%%%%%%%%%%%%%%%%%%%%%%%%%%%%%%%%%%%%%%%%%%%%%%%%%%%%%%%%%%

The following axiom system establishes the formal foundations of the welfare hierarchy.

% SRL-1
\begin{tcolorbox}[colback=gray!5!white, colframe=gray!75!black,
    title=Axiom SRL-1: Existence of Utility Bearer]
\textbf{Statement:}
\begin{equation}
\forall u \in \mathcal{U}: \exists! L \in \{0, \ldots, n\}, \; \exists! e \in \mathcal{E}_L: u = U_e^{(L)}
\end{equation}

\textbf{Interpretation:} Every utility value has exactly one bearer at exactly one level.
There is no ``free-floating'' utility without an entity to whom it belongs.

\textbf{Implication:} When measuring welfare, we must always specify \textit{whose} utility we mean.

\textbf{Testable Prediction:} Survey questions must identify the reference entity;
responses about ``general welfare'' without entity specification are meaningless.
\end{tcolorbox}

% SRL-2
\begin{tcolorbox}[colback=gray!5!white, colframe=gray!75!black,
    title=Axiom SRL-2: Level Boundedness]
\textbf{Statement:}
\begin{equation}
\forall L: U_e^{(L)} \in [U_{\min}^{(L)}, U_{\max}^{(L)}] \subset \mathbb{R}
\end{equation}

\textbf{Interpretation:} Utility at each level is bounded. There exists no infinite utility.

\textbf{Implication:} Enables meaningful aggregation; prevents utility monsters.

\textbf{Cross-Reference:} Connects to CORE-WHAT (C) for dimension-specific bounds.
\end{tcolorbox}

% SRL-3
\begin{tcolorbox}[colback=gray!5!white, colframe=gray!75!black,
    title=Axiom SRL-3: Non-Trivial Emergence]
\textbf{Statement:}
\begin{equation}
\exists L > 0: U_e^{(L)} \neq g\left(\{U_m^{(L-1)}\}_{m \in e}\right) \quad \forall g \in \mathcal{G}_{\text{additive}}
\end{equation}

\textbf{Interpretation:} For at least some higher levels, the collective utility is not merely
the sum (or any additive function) of constituent utilities. Genuine emergence exists.

\textbf{Implication:} Team spirit, organizational culture, and social trust are real phenomena,
not just aggregations of individual feelings.

\textbf{Testable Prediction:} Group performance measures will show variance unexplained
by individual member utilities (cf. team effectiveness research).
\end{tcolorbox}

% SRL-4
\begin{tcolorbox}[colback=gray!5!white, colframe=gray!75!black,
    title=Axiom SRL-4: Downward Causation]
\textbf{Statement:}
\begin{equation}
\frac{\partial U_i^{(0)}}{\partial U_e^{(L)}} \neq 0 \quad \text{for some } i \in e, \; L > 0
\end{equation}

\textbf{Interpretation:} Higher-level utility can causally affect lower-level utility.
A thriving organization can make individuals happier, not just vice versa.

\textbf{Implication:} Interventions at higher levels (organizational change, policy reform)
can improve individual welfare beyond what individual-level interventions achieve.

\textbf{Cross-Reference:} Connects to CORE-HOW (B) for the complementarity structure.
\end{tcolorbox}

% SRL-5
\begin{tcolorbox}[colback=gray!5!white, colframe=gray!75!black,
    title=Axiom SRL-5: Level-Specific Context]
\textbf{Statement:}
\begin{equation}
U_e^{(L)} = U_e^{(L)}(\Psi_L) \quad \text{where } \Psi_L \neq \Psi_{L'} \text{ for } L \neq L'
\end{equation}

\textbf{Interpretation:} Each level has its own relevant context dimensions.
The context affecting individual utility differs from that affecting organizational utility.

\textbf{Implication:} Context analysis (CORE-WHEN) must be level-specific.

\textbf{Cross-Reference:} Connects to CORE-WHEN (V) for $\Psi$-specification per level.
\end{tcolorbox}

% SRL-6
\begin{tcolorbox}[colback=gray!5!white, colframe=gray!75!black,
    title=Axiom SRL-6: Pareto Consistency]
\textbf{Statement:}
\begin{equation}
\left[ \forall L: U_L(x) \geq U_L(y) \right] \land \left[ \exists L: U_L(x) > U_L(y) \right] \Rightarrow W(x) > W(y)
\end{equation}

\textbf{Interpretation:} If option $x$ is at least as good as $y$ at every level, and strictly
better at some level, then total welfare under $x$ exceeds that under $y$.

\textbf{Implication:} The aggregation function $W$ respects Pareto improvements at all levels.

\textbf{Testable Prediction:} Policy choices that improve one level without harming others
should be unanimously preferred.
\end{tcolorbox}

% SRL-7
\begin{tcolorbox}[colback=gray!5!white, colframe=gray!75!black,
    title=Axiom SRL-7: No Double-Counting]
\textbf{Statement:}
\begin{equation}
W = \sum_{L=0}^{n} \omega_L \cdot U_L^{\text{net}} \quad \text{where } U_L^{\text{net}} = U_L - \sum_{L' < L} \pi_{L \leftarrow L'} \cdot U_{L'}
\end{equation}

\textbf{Interpretation:} When aggregating across levels, we must subtract the portion of
higher-level utility that is already counted at lower levels.

\textbf{Implication:} The ``inheritance coefficient'' $\pi_{L \leftarrow L'}$ captures how much
of level-$L'$ utility is embedded in level-$L$ utility.

\textbf{Testable Prediction:} Simple summation of utilities across levels will systematically
overestimate total welfare in nested structures.
\end{tcolorbox}

% SRL-8
\begin{tcolorbox}[colback=gray!5!white, colframe=gray!75!black,
    title=Axiom SRL-8: Weight Normalization]
\textbf{Statement:}
\begin{equation}
\sum_{L=0}^{n} \omega_L = 1 \quad \text{where } \omega_L \geq 0 \; \forall L
\end{equation}

\textbf{Interpretation:} Welfare weights across levels must sum to unity.
This is a normalization constraint ensuring well-defined total welfare.

\textbf{Implication:} The choice of $\{\omega_L\}$ is a normative decision about whose
interests count and how much.

\textbf{Cross-Reference:} Parameter values for $\omega_L$ come from CORE-WHERE (BBB).
\end{tcolorbox}


%%%%%%%%%%%%%%%%%%%%%%%%%%%%%%%%%%%%%%%%%%%%%%%%%%%%%%%%%%%%%%%%%%%%%%%%%%%%%%%
\section{Aggregation Theory}
\label{sec:who-aggregation}
%%%%%%%%%%%%%%%%%%%%%%%%%%%%%%%%%%%%%%%%%%%%%%%%%%%%%%%%%%%%%%%%%%%%%%%%%%%%%%%

\subsection{The General Aggregation Formula}

Combining axioms SRL-6, SRL-7, and SRL-8, the total welfare function is:

\begin{equation}
W = \sum_{L=0}^{n} \omega_L \cdot \left[ U_L - \sum_{L' < L} \pi_{L \leftarrow L'} \cdot U_{L'} \right]
\label{eq:who-full-aggregation}
\end{equation}

\subsection{Special Cases}

\paragraph{Pure Individualism ($\omega_0 = 1$):}
When only individual utility matters:
\[
W = U_0 = \sum_{i=1}^{N} u_i
\]
This recovers classical utilitarian welfare.

\paragraph{Pure Collectivism ($\omega_n = 1$):}
When only societal utility matters:
\[
W = U_n^{\text{net}} = U_n - \sum_{L < n} \pi_{n \leftarrow L} \cdot U_L
\]

\paragraph{Balanced Weighting:}
The empirically calibrated case with $\omega_L > 0$ for multiple levels,
reflecting that both individual and collective welfare matter.

\subsection{Determining Welfare Weights}

The welfare weights $\{\omega_L\}$ are normative parameters requiring:

\begin{enumerate}
\item \textbf{Ethical justification}: What theory supports these weights?
\item \textbf{Democratic legitimacy}: Do citizens endorse these weights?
\item \textbf{Context sensitivity}: Should weights vary by decision domain?
\end{enumerate}

See Appendix BBB (CORE-WHERE) for calibration methods.


%%%%%%%%%%%%%%%%%%%%%%%%%%%%%%%%%%%%%%%%%%%%%%%%%%%%%%%%%%%%%%%%%%%%%%%%%%%%%%%
\section{Cross-Level Externalities}
\label{sec:who-externalities}
%%%%%%%%%%%%%%%%%%%%%%%%%%%%%%%%%%%%%%%%%%%%%%%%%%%%%%%%%%%%%%%%%%%%%%%%%%%%%%%

\subsection{Definition}

\begin{definition}[Cross-Level Externality]
A \textbf{cross-level externality} exists when the utility at level $L$ directly affects
utility at level $L' \neq L$:
\begin{equation}
\frac{\partial U_{L'}}{\partial U_L} \neq 0
\end{equation}
\end{definition}

\subsection{Types of Cross-Level Externalities}

\begin{table}[htbp]
\centering
\caption{Cross-Level Externality Types}
\label{tab:who-externalities}
\renewcommand{\arraystretch}{1.3}
\begin{tabular}{@{}lll@{}}
\toprule
\textbf{Direction} & \textbf{Type} & \textbf{Example} \\
\midrule
$L \to L+1$ (Upward) & Constitutive & Individual efforts create team success \\
$L+1 \to L$ (Downward) & Enabling & Organizational culture enables individual flourishing \\
$L \to L$ (Horizontal) & Peer Effects & One team's innovation inspires another \\
$L \to L+k$ (Skip-Level) & Structural & Individual voting shapes societal outcomes \\
\bottomrule
\end{tabular}
\end{table}


%%%%%%%%%%%%%%%%%%%%%%%%%%%%%%%%%%%%%%%%%%%%%%%%%%%%%%%%%%%%%%%%%%%%%%%%%%%%%%%
\section{Integration with Other CORE Appendices}
\label{sec:who-integration}
%%%%%%%%%%%%%%%%%%%%%%%%%%%%%%%%%%%%%%%%%%%%%%%%%%%%%%%%%%%%%%%%%%%%%%%%%%%%%%%

\subsection{Connection to CORE-WHAT (Appendix C)}

\begin{tcolorbox}[colback=yellow!5!white, colframe=yellow!75!black,
    title=Integration: CORE-WHO $\leftrightarrow$ CORE-WHAT]
\textbf{What WHO provides to WHAT:}
\begin{itemize}[nosep]
\item The bearer specification: \textit{whose} utility dimensions?
\item Level-specific dimension relevance: different $d \in \text{FEPSDE}$ matter at different $L$
\end{itemize}

\textbf{What WHO requires from WHAT:}
\begin{itemize}[nosep]
\item Content for $U_L$: what are the dimensions of utility at each level?
\item Comparability across levels: can we compare $U_L$ and $U_{L'}$?
\end{itemize}

\textbf{Mathematical Interface:}
\begin{equation}
U_e^{(L)} = \sum_{d \in \text{FEPSDE}} \alpha_d^{(L)} \cdot u_d^{(L)}
\end{equation}
where $\alpha_d^{(L)}$ weights dimensions by level.
\end{tcolorbox}

\subsection{Connection to CORE-HOW (Appendix B)}

\begin{tcolorbox}[colback=yellow!5!white, colframe=yellow!75!black,
    title=Integration: CORE-WHO $\leftrightarrow$ CORE-HOW]
\textbf{What WHO provides to HOW:}
\begin{itemize}[nosep]
\item The level structure for complementarity: $\gamma_{LL'}$ between levels
\item Entities whose interactions to model
\end{itemize}

\textbf{What WHO requires from HOW:}
\begin{itemize}[nosep]
\item Cross-level complementarity values $\gamma_{L \leftrightarrow L'}$
\item Whether levels are complements ($\gamma > 0$) or substitutes ($\gamma < 0$)
\end{itemize}

\textbf{Mathematical Interface:}
\begin{equation}
\frac{\partial^2 W}{\partial U_L \partial U_{L'}} = \gamma_{LL'}
\end{equation}
\end{tcolorbox}

\subsection{Connection to CORE-WHEN (Appendix V)}

\begin{tcolorbox}[colback=yellow!5!white, colframe=yellow!75!black,
    title=Integration: CORE-WHO $\leftrightarrow$ CORE-WHEN]
\textbf{What WHO provides to WHEN:}
\begin{itemize}[nosep]
\item Level-specific context relevance: which $\Psi$ dimensions matter at each $L$
\item Entities whose context to analyze
\end{itemize}

\textbf{What WHO requires from WHEN:}
\begin{itemize}[nosep]
\item Context-dependent welfare weights: $\omega_L(\Psi)$
\item When aggregation rules should change
\end{itemize}

\textbf{Mathematical Interface:}
\begin{equation}
\omega_L = \omega_L(\Psi) \quad \text{where } \Psi = (\Psi_I, \Psi_S, \Psi_C, \Psi_K, \Psi_E, \Psi_T, \Psi_M, \Psi_F)
\end{equation}
\end{tcolorbox}


%%%%%%%%%%%%%%%%%%%%%%%%%%%%%%%%%%%%%%%%%%%%%%%%%%%%%%%%%%%%%%%%%%%%%%%%%%%%%%%
\section{Worked Example}
\label{sec:who-example}
%%%%%%%%%%%%%%%%%%%%%%%%%%%%%%%%%%%%%%%%%%%%%%%%%%%%%%%%%%%%%%%%%%%%%%%%%%%%%%%

\subsection{Example: Organizational Restructuring Decision}

\paragraph{Setup.}
A company considers restructuring that would:
\begin{itemize}
\item Increase $U_0$ (individual utility) for 70\% of employees by CHF 5,000/year
\item Decrease $U_0$ for 30\% of employees (layoffs) by CHF 50,000/year
\item Increase $U_3$ (organizational utility) by CHF 2M through improved efficiency
\item Decrease $U_2$ (team utility) in affected departments
\end{itemize}

\paragraph{Application.}
Using the aggregation formula with Swiss context weights from BBB:
\begin{align}
W &= \omega_0 \cdot U_0^{\text{net}} + \omega_2 \cdot U_2^{\text{net}} + \omega_3 \cdot U_3^{\text{net}} \\
  &= 0.5 \cdot [0.7 \times 5000 - 0.3 \times 50000] \\
  &\quad + 0.2 \cdot [\Delta U_2] \\
  &\quad + 0.3 \cdot [2,000,000 - \pi_{3 \leftarrow 0} \cdot \Delta U_0]
\end{align}

\paragraph{Result.}
The multi-level analysis reveals that organizational gains may be offset by
individual and team-level losses. A pure organizational perspective ($\omega_3 = 1$)
would overestimate net benefit.

\paragraph{Discussion.}
This example illustrates why single-level analysis is insufficient:
different stakeholders experience the same decision differently,
and the welfare hierarchy makes these trade-offs explicit.


%%%%%%%%%%%%%%%%%%%%%%%%%%%%%%%%%%%%%%%%%%%%%%%%%%%%%%%%%%%%%%%%%%%%%%%%%%%%%%%
\section{Implications for Practice}
\label{sec:who-practice}
%%%%%%%%%%%%%%%%%%%%%%%%%%%%%%%%%%%%%%%%%%%%%%%%%%%%%%%%%%%%%%%%%%%%%%%%%%%%%%%

\begin{tcolorbox}[colback=green!5!white, colframe=green!75!black,
    title=Practical Implications]
\begin{enumerate}
\item \textbf{Always Specify the Level:} When measuring ``welfare,'' explicitly state
      at which level(s) and for which entities.
\item \textbf{Beware Simple Aggregation:} Summing utilities across levels without
      accounting for nesting ($\pi_{L \leftarrow L'}$) leads to double-counting.
\item \textbf{Make Weights Explicit:} Policy debates often mask implicit disagreements
      about $\omega_L$ --- surface these for transparent decision-making.
\item \textbf{Consider Emergence:} Some collective goods (trust, culture) exist at higher
      levels and cannot be achieved through individual-level interventions alone.
\end{enumerate}
\end{tcolorbox}


%%%%%%%%%%%%%%%%%%%%%%%%%%%%%%%%%%%%%%%%%%%%%%%%%%%%%%%%%%%%%%%%%%%%%%%%%%%%%%%
\section{Summary}
\label{sec:who-summary}
%%%%%%%%%%%%%%%%%%%%%%%%%%%%%%%%%%%%%%%%%%%%%%%%%%%%%%%%%%%%%%%%%%%%%%%%%%%%%%%

\begin{tcolorbox}[colback=gray!5!white, colframe=gray!75!black,
    title=Appendix AAA Summary: CORE-WHO]

\textbf{Core Contribution:}
Establishes that utility exists at multiple hierarchical levels $L \in \{0, \ldots, n\}$,
from individual to society, with formal rules for cross-level aggregation.

\textbf{Key Formula:}
\begin{equation}
W = \sum_{L=0}^{n} \omega_L \cdot \left[ U_L - \sum_{L' < L} \pi_{L \leftarrow L'} \cdot U_{L'} \right]
\end{equation}

\textbf{Key Insights:}
\begin{itemize}[nosep]
\item WHO must be answered before WHAT, HOW, or WHEN
\item Higher-level utilities are not mere aggregations (SRL-3: emergence)
\item Cross-level externalities require explicit modeling (SRL-4: downward causation)
\item Welfare weights $\omega_L$ are normative choices that should be transparent
\end{itemize}

\textbf{Next Steps:}
See Appendix C (CORE-WHAT) for the content of utility at each level.
See Appendix B (CORE-HOW) for complementarity across levels.

\end{tcolorbox}


%%%%%%%%%%%%%%%%%%%%%%%%%%%%%%%%%%%%%%%%%%%%%%%%%%%%%%%%%%%%%%%%%%%%%%%%%%%%%%%
\section{Glossary of Symbols}
\label{sec:who-glossary}
%%%%%%%%%%%%%%%%%%%%%%%%%%%%%%%%%%%%%%%%%%%%%%%%%%%%%%%%%%%%%%%%%%%%%%%%%%%%%%%

\begin{tcolorbox}[colback=blue!5!white, colframe=blue!75!black]
\textbf{Central Reference:} For complete symbol definitions, see
\textbf{Appendix GLS} (\textit{Glossary and Notation}).

This section lists only symbols introduced or specialized in this appendix.
\end{tcolorbox}

\begin{table}[htbp]
\centering
\caption{Symbols Introduced in Appendix AAA}
\label{tab:who-symbols}
\renewcommand{\arraystretch}{1.3}
\begin{tabular}{@{}clp{6cm}c@{}}
\toprule
\textbf{Symbol} & \textbf{Name} & \textbf{Definition} & \textbf{Status} \\
\midrule
$L$ & Level & Aggregation tier (0 = individual, $n$ = society) & NEW \\
$\mathcal{E}_L$ & Entity Set & Set of utility-bearing entities at level $L$ & NEW \\
$U_e^{(L)}$ & Entity Utility & Utility of entity $e$ at level $L$ & NEW \\
$\omega_L$ & Welfare Weight & Normative weight for level $L$ in aggregation & NEW \\
$\pi_{L \leftarrow L'}$ & Inheritance Coef. & Share of $L'$-utility embedded in $L$-utility & NEW \\
$\Xi_L$ & Emergence Term & Non-reducible component of level-$L$ utility & NEW \\
$W$ & Total Welfare & Aggregated welfare across all levels & \checkmark \\
\bottomrule
\end{tabular}
\end{table}


%%%%%%%%%%%%%%%%%%%%%%%%%%%%%%%%%%%%%%%%%%%%%%%%%%%%%%%%%%%%%%%%%%%%%%%%%%%%%%%
\section{Critical Foundations and Objections}
\label{sec:who-critical}
%%%%%%%%%%%%%%%%%%%%%%%%%%%%%%%%%%%%%%%%%%%%%%%%%%%%%%%%%%%%%%%%%%%%%%%%%%%%%%%

\subsection{Objection 1: Methodological Individualism}

\begin{tcolorbox}[colback=red!5!white, colframe=red!50!black,
    title=Objection: Only Individuals Have Utility]
\textit{``Collectives are not agents. Only individuals feel pleasure and pain.
Speaking of `organizational utility' or `societal welfare' as distinct from
individual welfare is a category mistake.''}
\end{tcolorbox}

\textbf{Response:}
The objection conflates \textit{phenomenal experience} with \textit{welfare}.
While only individuals \textit{experience}, there exist collective states
(trust, culture, social capital) that:
\begin{enumerate}
\item Cannot be decomposed into individual states
\item Causally affect individual welfare (downward causation)
\item Are legitimate targets of policy
\end{enumerate}

\textbf{Empirical Evidence:}
\citet{putnam2000bowling} shows social capital has effects not reducible
to individual characteristics. Team effectiveness research consistently
finds group-level variance unexplained by member composition.

\subsection{Objection 2: Double-Counting Problem}

\begin{tcolorbox}[colback=red!5!white, colframe=red!50!black,
    title=Objection: Multi-Level Aggregation Double-Counts]
\textit{``If we count both individual utility and group utility, we're counting
the same things twice, since groups are made of individuals.''}
\end{tcolorbox}

\textbf{Response:}
This is why Axiom SRL-7 introduces the inheritance coefficient $\pi_{L \leftarrow L'}$.
Proper multi-level aggregation subtracts the portion already counted at lower levels.
The formula in Equation~\ref{eq:who-full-aggregation} explicitly avoids double-counting.

\subsection{Objection 3: Interpersonal Comparison Problem}

\begin{tcolorbox}[colback=red!5!white, colframe=red!50!black,
    title=Objection: Cross-Level Comparison is Meaningless]
\textit{``How can we compare individual utility with organizational utility?
They're measured in different units and have no common scale.''}
\end{tcolorbox}

\textbf{Response:}
This objection applies equally to within-level comparisons between individuals.
The framework does not require cardinal comparability; it requires only:
\begin{enumerate}
\item Ordinal consistency within each level
\item Normative weights $\omega_L$ that express trade-off preferences
\end{enumerate}

See Appendix C (CORE-WHAT) for dimensionless utility normalization.


%%%%%%%%%%%%%%%%%%%%%%%%%%%%%%%%%%%%%%%%%%%%%%%%%%%%%%%%%%%%%%%%%%%%%%%%%%%%%%%
\section{Open Issues and Research Agenda}
\label{sec:who-open}
%%%%%%%%%%%%%%%%%%%%%%%%%%%%%%%%%%%%%%%%%%%%%%%%%%%%%%%%%%%%%%%%%%%%%%%%%%%%%%%

\begin{table}[htbp]
\centering
\caption{Research Roadmap for CORE-WHO}
\label{tab:who-roadmap}
\renewcommand{\arraystretch}{1.3}
\begin{tabular}{@{}clccl@{}}
\toprule
\textbf{\#} & \textbf{Issue} & \textbf{Priority} & \textbf{Status} & \textbf{Requirements} \\
\midrule
1 & Calibrate $\omega_L$ cross-culturally & HIGH & Open & Survey data from 20+ countries \\
2 & Estimate $\pi_{L \leftarrow L'}$ empirically & HIGH & Open & Multi-level panel data \\
3 & Formalize temporal discounting by level & MEDIUM & Open & Intertemporal choice theory \\
4 & Non-human entities (AI, ecosystems) & MEDIUM & Open & Philosophical groundwork \\
5 & Democratic weight determination & LOW & Open & Deliberative democracy methods \\
\bottomrule
\end{tabular}
\end{table}


%%%%%%%%%%%%%%%%%%%%%%%%%%%%%%%%%%%%%%%%%%%%%%%%%%%%%%%%%%%%%%%%%%%%%%%%%%%%%%%
\section{References}
\label{sec:who-references}
%%%%%%%%%%%%%%%%%%%%%%%%%%%%%%%%%%%%%%%%%%%%%%%%%%%%%%%%%%%%%%%%%%%%%%%%%%%%%%%

\begin{tcolorbox}[colback=blue!5!white, colframe=blue!75!black]
\textbf{Master Bibliography:} For complete reference details, see
\texttt{bcm\_master.bib} in the repository.
\end{tcolorbox}

\subsection{Key References for This Appendix}

\begin{itemize}[nosep]
\item \textbf{Social Welfare Functions:} \citet{sen1970collective}, \citet{arrow1951social}
\item \textbf{Multi-Level Analysis:} \citet{coleman1990foundations}, \citet{raudenbush2002hierarchical}
\item \textbf{Social Capital:} \citet{putnam2000bowling}, \citet{coleman1988social}
\item \textbf{Organizational Behavior:} \citet{kozlowski2000multilevel}
\item \textbf{Team Effectiveness:} \citet{hackman2002groups}
\end{itemize}

% Include master bibliography
\nocite{bcm_master}

%%%%%%%%%%%%%%%%%%%%%%%%%%%%%%%%%%%%%%%%%%%%%%%%%%%%%%%%%%%%%%%%%%%%%%%%%%%%%%%
% CROSS-REFERENCE FOOTER
%%%%%%%%%%%%%%%%%%%%%%%%%%%%%%%%%%%%%%%%%%%%%%%%%%%%%%%%%%%%%%%%%%%%%%%%%%%%%%%

\vspace{1em}
\hrule
\vspace{0.5em}

\noindent\textit{Cross-references:}
Appendix GLS (Glossary),
Appendix C (CORE-WHAT),
Appendix B (CORE-HOW),
Appendix V (CORE-WHEN),
Appendix BBB (CORE-WHERE),
Appendix AU (CORE-AWARE),
Appendix CAL (METHOD-LLMMC).

\noindent\textit{Master Bibliography:} \texttt{bcm\_master.bib}

% =============================================================================
% END OF APPENDIX AAA
% =============================================================================
